% concatenated from main.tex
\documentclass[11pt]{article}

\usepackage{amsmath, amssymb, amsthm, mathtools}
\usepackage[margin=1in]{geometry}
\usepackage{enumitem}
\usepackage{hyperref}
\usepackage{cleveref}
\usepackage{mleftright}
\usepackage{xcolor}
\hypersetup{
	colorlinks=true,
	linkcolor=blue!50!black, % Deep blue for internal links (sections, equations)
	citecolor=green!50!black, % Deep green for citations
	urlcolor=red!50!black % Deep red/maroon for web URLs
}
\usepackage{algorithm}
\usepackage{algorithmicx}
\usepackage{algpseudocode}

\newtheorem{theorem}{Theorem}[section]
\newtheorem{lemma}[theorem]{Lemma}
\newtheorem{corollary}[theorem]{Corollary}

\theoremstyle{remark}
\newtheorem{remark}[theorem]{Remark}

\theoremstyle{definition}
\newtheorem{definition}{Definition}


\newcommand{\Sph}{\mathbb{S}}
\newcommand{\R}{\mathbb{R}}
\newcommand{\Z}{\mathbb{Z}}
\newcommand{\C}{\mathbb{C}}
\newcommand{\eps}{\epsilon}
\newcommand{\cP}{\mathcal{P}}
\newcommand{\cA}{\mathcal{A}}
\newcommand{\cB}{\mathcal{B}}
\newcommand{\cC}{\mathcal{C}}
\newcommand{\cD}{\mathcal{D}}
\newcommand{\cE}{\mathcal{E}}
\newcommand{\cF}{\mathcal{F}}
\newcommand{\cG}{\mathcal{G}}
\newcommand{\cM}{\mathcal{M}}
\newcommand{\cN}{\mathcal{N}}
\newcommand{\cR}{\mathcal{R}}
\newcommand{\cS}{\mathcal{S}}
\newcommand{\ip}[2]{\langle #1,#2\rangle}
\newcommand{\norm}[1]{\left\|#1\right\|}
\newcommand{\embeds}[1]{\overset{#1}{\hookrightarrow}}
\newcommand{\op}{\textrm{op}}
\DeclareMathOperator{\supp}{supp}
\DeclareMathOperator{\poly}{poly}
\DeclareMathOperator{\polylog}{polylog}
\DeclareMathOperator{\range}{range}
\DeclareMathOperator{\diam}{diam}
\DeclareMathOperator*{\E}{\mathbb{E}}
\DeclareMathOperator*{\Var}{Var}
\DeclareMathOperator{\sgn}{sgn}
\DeclareMathOperator{\diag}{diag}
\DeclareMathOperator{\rk}{rank}
\DeclareMathOperator{\spn}{span}
\newcommand{\TVnorm}[1]{\norm{#1}_{\mathrm{TV}}}

\title{A Near-Optimal Lower Bound for $\ell_p$-Subspace Embeddings, $1\leq p<2$}
\author{Yi Li \\ Nanyang Technological University \\ \texttt{yili@ntu.edu.sg}}
\date{}
\begin{document}
\maketitle

\begin{abstract}
	For $d \geq 2$, $p \geq 1$ and $\epsilon > 0$, let $N_p(d,\epsilon)$ be the smallest integer $N$ such that for every integer $n$ and every $A\in\mathbb{R}^{n\times d}$, there exists a matrix $\Phi\in\mathbb{R}^{N\times n}$ satisfying
	$(1-\epsilon)\lVert Ax\rVert_p\leq \lVert\Phi A x\rVert_p\leq (1+\epsilon)\lVert Ax\rVert_p$ for all $x\in\mathbb{R}^d$.
	For every constant $p\geq 1$ with $p\not\in 2\mathbb{Z}$, when $d\gtrsim_p \log(1/\epsilon)$, the bound
	\[
		N_p(d,\epsilon) \gtrsim_{p} \frac{d}{\epsilon^2 \operatorname{polylog}(d/\epsilon)}
	\]
	is established. This improves the previous lower bound $\Omega(1/(\epsilon^2\operatorname{polylog}(1/\epsilon)))$ due to Li et al. (SICOMP 2021) and is optimal up to logarithmic factors for $1\leq p<2$.  The central technical idea originated from ChatGPT 5.6 Sol.
\end{abstract}

\section{Introduction}

The $\ell_p$-subspace embedding problem is a classical problem in the local theory of Banach spaces. In its original form, it considers embedding a finite-dimensional subspace of $L^p[0,1]$. Since such a subspace can be embedded with small distortion into $\ell_p^n$ with a sufficiently large $n$, we present the $\ell_p$-subspace embedding problem in the following form.

\begin{definition}[$\ell_p$-subspace embedding]
	For $d\geq 1$, $p\geq 1$ and $\eps>0$, let $N_p(d,\eps)$ be the smallest integer $N$ such that for every integer $n$ and every $A\in\R^{n\times d}$, there exists a matrix $\Phi \in \R^{N\times n}$ satisfying
	\[
		(1-\epsilon)\norm{Ax}_{p}\leq \norm{\Phi A x}_p \leq (1+\eps)\norm{Ax}_p,	\qquad \forall x\in \R^d.
	\]
\end{definition}

The problem has a long research history, dating back to the 1980s; see a comprehensive survey in~\cite{handbook:19}. Different results in the literature obtain similar upper bounds for $N_p(d,\eps)$ of the form
\begin{equation}\label{eqn:N_p(d,eps)}
	N_p(d,\eps) \leq C \eps^{-2}d^{\max\{1, p/2\}}\log^c (d/\eps)
\end{equation}
but with different logarithmic factors. Here, $C,c > 0$ are absolute constants. The proofs can be found in, e.g., the monograph by Ledoux and Talagrand~\cite{LT91}. It is known that when $p\in 2\Z$, $N_p(d,\eps) \lesssim p^{-2}\eps^{-2}\binom{d+p/2-1}{p/2}$ for $\eps\leq 1/p$, without additional logarithmic factors~\cite{schechtman:tight}. Very recently, Reis and Rothvoss~\cite{RR26} obtained $N_1(d,\eps) \lesssim d/\eps^2$, removing the logarithmic factor for $p = 1$.

For fixed $\eps$, the dependence $d^{\max\{1,p/2\}}$ in \eqref{eqn:N_p(d,eps)} has long been known to be tight. Regarding the dependence on $\eps$, it was proved a few years ago that $\eps^{-2}$ is tight up to logarithmic factors when $d\gtrsim \log(1/\eps)$ and $p\not\in 2\Z$~\cite{LWW21}. Here, some constraint on $d$ is necessary, since it is known that one can do strictly better than $\eps^{-2}$ when $d$ is a constant. In particular, for fixed $d$ and $p\not\in 2\Z$, it is now known that $N_p(d,\eps) \asymp_{d,p} (1/\eps)^{2(d-1)/(d+2p)}$. The lower bound was proved in~\cite{LLW23} and the upper bound was proved very recently in~\cite{L26:const_dim}. The assumption $p\not\in 2\Z$ is also necessary, since isometric embeddings are known to exist for $p\in 2\Z$~\cite{handbook:21}, and so there cannot be a lower bound involving $\eps$.

For $p\in[1,2)$ and $d\gtrsim \log(1/\eps)$, Li et al.~\cite{LWW21} also proved that the factor $\eps^{-2}d$ is tight up to logarithmic factors when the embedding matrix $\Phi$ is restricted to be a row-sampling matrix (with rescaling).
However, it was not known in general whether the combined dependence $\eps^{-2}d^{\max\{1,p/2\}}$ is tight. Our main result gives a lower bound of $\widetilde{\Omega}_p(d/\eps^2)$ for every constant $p\geq 1$ with $p\notin 2\mathbb{Z}$. In particular, for $1\leq p<2$, this matches the known upper bound up to logarithmic factors.

The proof follows the framework of Li et al.~\cite{LWW21}. They introduced a data structure problem called the $\ell_p$-subspace sketch problem, proved a bit lower bound for it, and then converted this bit lower bound to a dimension lower bound. We first define the following for-all version of the $\ell_p$-subspace sketch problem.
\begin{definition}[$\ell_p$-subspace sketch, for-all version]
	Given an $n\times d$ matrix $A$ with entries specified by $O(\log(nd))$ bits,
	an accuracy parameter $\eps  > 0$ and a constant $p\geq 1$, design a data structure
	$Q_p$ so that with probability at least 0.9, $Q_p(x) = (1 \pm  \eps)\|Ax\|_p^p$ for all $x \in \R^d$.
\end{definition}

Throughout the paper, let $\alpha_p=\max\{1,p/2\}$. We have the following lower bound for the $\ell_p$-subspace sketch problem.

\begin{theorem}\label{thm:subspace_sketch_lower_bound}
	Let $p \in  [1, \infty) \setminus 
	2\mathbb{Z}$ be a constant. For any
	$d = \Omega_p(\log(1/\eps))$ and
	$n = \widetilde\Omega_p(d^{\alpha_p}/\eps^{2+2\alpha_p})$, we have that $\Omega(d^2/(\eps^2\polylog(1/\eps)))$ bits are necessary to solve the $\ell_p$-subspace sketch problem.
\end{theorem}

Following the same approach as in \cite{LWW21} to convert a bit lower bound to a dimension lower bound, we obtain the following theorem.
\begin{theorem}\label{thm:dimension_lower_bound}
	Let $p\in [1, \infty) \setminus 
	2\mathbb{Z}$ be a constant.
	For $d=\Omega_p(\log(1/\eps))$, it holds that
	\[
		N_p(d,\epsilon) \gtrsim_{p} \frac{d}{\eps^2 \polylog(d/\eps)}.
	\]
\end{theorem}
When $p\in [1,2)$, this lower bound is tight up to logarithmic factors and the $\ell_p$-subspace embedding problem is almost settled. For $p>2$, proving a lower bound of $\widetilde{\Omega}(d^{p/2}/\eps^2)$ for $N_p(d,\eps)$ remains open.

\subsection{Review of the previous construction}
Before proving Theorem~\ref{thm:subspace_sketch_lower_bound}, we first recall the strategy underlying the earlier $\tilde{\Omega}(d/\eps^2)$ lower bound of \cite{LWW21} for the for-each $\ell_p$-subspace sketch problem. For easier comparison, we shall describe the construction for a for-all version.

The argument begins with a basic hard instance $A\in \R^{2^k\times k}$, where $k$ is even and $2^k=\widetilde{\Theta}_p(\eps^{-2})$, for which any $\ell_p$-subspace sketch requires $2^k/\poly(k) = \tilde{\Omega}(\eps^{-2})$ bits. To extend the construction to general dimension $d$, one places independent copies of $A$ along the diagonal of a block-diagonal matrix with $d$ columns. We therefore review the lower-bound argument for this basic instance.

Let the index universe $U = \{-1,1\}^k$. Let $M\in\R^{2^k\times 2^k}$ with $M_{i,j} = |\langle i,j\rangle|^p$ for $i,j\in U$.  The matrix $M$ has $r = \binom{k}{k/2}$ eigenvalues $\lambda$ satisfying $|\lambda| \gtrsim_p \sqrt{2^k/k}$. The normalized Hadamard matrix $H$ diagonalizes $M$; after permuting the columns of $H$, assume that the first $r$ diagonal entries of
$H^\top MH$ are all $\lambda$. Define the spectrum-truncated version of $M$ as
\[
	\widetilde{M} = H \diag\{\underbrace{\lambda,\dots,\lambda}_{r\text{ copies}},0,\dots,0\} H^\top.
\]
One then selects a subset $I\subseteq U$ with $|I| = \Omega(r/\log r)$ and decomposes $\widetilde{M}_i = R_i + P_i$ for each $i\in I$ such that each $R_i\in\range(\widetilde{M})$, $\{R_i\}_{i\in I}$ are orthogonal, $\norm{R_i}_2\geq 0.99\|\widetilde{M}_i\|_2$, and $P_i$ is a small projection onto the span of the other rows $R_j$ ($j\neq i$).
The message encoded in the hard instance is the sign vector $(s_i)_{i\in I}\in\{-1,1\}^I$. Given this sign vector, define a vector 
\begin{equation}\label{eqn:old_x}
	x = \sum_{i\in I} s_i \frac{R_i}{\norm{R_i}_2}.
\end{equation}
For an appropriately chosen $\Delta$, with constant probability over uniformly random signs $s_i$, both $x_j+\Delta\geq 0$ for all $j\in U$ and
$|\langle P_i,x\rangle|\leq (1/10)\norm{R_i}_2$ for all $i\in I$ hold simultaneously. For every such sign vector, define weights $y_j=(x_j+\Delta)^{1/p}$ and rows $A_j=y_j\cdot j^\top$ for $j\in U$.

For a query vector $i \in I$, $(Ai)_j = y_j \ip{j}{i}$, and hence 
\[
	\norm{Ai}_p^p = \sum_j (x_j + \Delta)|\ip{j}{i}|^p = (Mx)_i + \Delta(M\mathbf{1})_i.
\]
The second term $\Delta(M\mathbf{1})_i$ is known exactly. Since $x\in\range(\widetilde{M})$ and $M$ and $\widetilde{M}$ agree on this subspace, it holds that $Mx = \widetilde{M}x$ and so
\begin{equation}\label{eqn:(Mx)_i_old}
	(Mx)_i = s_i \norm{R_i}_2 + \langle P_i, x\rangle.
\end{equation}
Hence there is a set $\mathcal{S}\subseteq\{-1,1\}^{I}$ of size $2^{\Omega(|I|)}$ for which both properties hold. Consequently, for every $(s_i)_{i\in I}\in\mathcal S$, an $(1\pm\eps)$-approximation $Q_p(i)$ to $\norm{Ai}_p^p$ allows one to recover $(Mx)_i$ approximately, and so one can recover $s_i$ from
\[
	s_i = \sgn(Q_p(i) - \Delta(M\mathbf{1})_i).
\]
Distinct sign vectors in $\mathcal{S}$ therefore require distinct contents stored in the data structure. Hence the data structure must have at least
\[
	\Omega(|I|) = \Omega\mleft(\frac{r}{\log r}\mright) = \Omega\mleft(\frac{2^k}{\poly(k)}\mright) = \Omega\mleft(\frac{1}{\eps^2\polylog(1/\eps)}\mright)
\]
bits.

\subsection{Our modification}
The starting point of our construction is the analogous spectral construction for a related matrix $M\in\R^{2^k\times 2^k}$, with $2^k = \widetilde{\Theta}_p(1/\eps^2)$. The key change is that each scalar sign $s_i$ is replaced by an $s\times s$ sign matrix $S_i$, where $s=(d-k)/2$. Thus, instead of encoding one bit at each index $i\in I$, the hard instance encodes $\Theta(s^2)$ bits at each such index.

To make these matrix-valued signs visible to norm queries, we introduce
\[
	C_j
	=
	I_{2s}
	+
	\eta
	\begin{pmatrix}
		0 & X_j\\
		X_j^\top & 0
	\end{pmatrix},
	\qquad\text{where}\quad
	X_j
	=
	\sum_{i\in I}\frac{(R_i)_j}{\norm{R_i}_2}\frac{S_i}{\sqrt{s}}.
\]
Note that the family $\{X_j\}_{j\in U}$ is the matrix-valued analogue of the vector $x$ in \eqref{eqn:old_x} for the previous construction.

The hard matrix $A$ is built so that its $\ell_p$-norms approximate the analytic quantity
\[
	F_S(v,\tau,w)
	=
	\sum_{j\in U}
	\mleft(\ip{j}{v}^2+\tau w^\top C_jw\mright)^{p/2}.
\]
More precisely, for query vectors of the form $x=(v,\sqrt{\tau}w)\in\R^k\oplus \R^{2s}$,
\[
	\norm{Ax}_p^p = (1\pm\eps)F_S(v,\tau,w).
\]
Thus the sketch gives approximate values of $F_S$. By evaluating these values at a number of carefully chosen values of $\tau$, one can approximate the derivative $\partial_\tau F_S(i,0,w)$.

For $v=i\in I$, this derivative takes the form
\[
	\partial_\tau F_S(i,0,w)
	=
	\frac p2 b_0
	+
	\frac p2 \eta\, w^\top
	\begin{pmatrix}
		0 & Y_i\\
		Y_i^\top & 0
	\end{pmatrix}
	w,
\]
where $b_0$ is independent of $i$ and can be computed exactly, and
\[
	Y_i
	=
	\norm{R_i}_2\frac{S_i}{\sqrt{s}}+E_i
\]
is the analogue of $(Mx)_i$ in \eqref{eqn:(Mx)_i_old}, where $E_i$ is a small error term. Varying $w$ gives access to the quadratic form of the block matrix above. Using that $E_i$ is small, one can show that accurate estimates of these quadratic forms for all $i\in I$ and $w\in\Sph^{2s-1}$ distinguish the members of a family of $2^{\Omega(|I|s^2)}$ sign-matrix collections. Since $|I|=\widetilde{\Omega}_p(\eps^{-2})$ and $s=\Theta(d)$, this gives the desired $\widetilde{\Omega}_p(d^2/\eps^2)$ bit lower bound.

\section{Hard Instance}
As previewed, we shall replace signs $s_i$ in the construction of \cite{LWW21} with sign matrices $S_i$ and aim to recover the sign matrices $S_i$. Fix a sufficiently small absolute constant $\theta>0$. Specifically, let $L\asymp \log(1/\eps)$ and choose $k=4h+1$ to be the largest integer of this form satisfying, for a sufficiently small constant $c_{p,\theta}>0$ depending only on $p$ and $\theta$,
\begin{equation}\label{eqn:assumption_k}
	2^k \leq \frac{c_{p,\theta}}{\eps^2 L^4 k^{p+3/2}}.
\end{equation}
Let $s = \lfloor (d-k)/2\rfloor$. Our query space is $\R^{k+2s} = \R^k \oplus \R^{2s}$. If $k + 2s < d$, we can append a zero column. To simplify the notation, we assume that $d = k+2s$ below. By adjusting the constants, the assumption that $d\gtrsim_p \log(1/\eps)$ implies that $s = \Theta(d)$.

Let $M\in\R^{2^k\times 2^k}$ be defined as $M_{i,j} = |\ip{i}{j}|^{p-2}$ for $i,j\in U$. Since $k$ is odd, $\ip{i}{j}\neq 0$ for all $i,j\in U$, so $M$ is well defined. Let $m = (k-1)/2$. 	Similarly to \cite{LWW21}, we have the following property of $M$.

\begin{lemma}\label{lem:eigenvalue}
Let $p\in [1,\infty)\setminus 2\mathbb{Z}$. The matrix $M$ has
\[
	r = \binom{k}{m}
\]
eigenvalues equal to $\lambda_{k,p}$, where
\[
	|\lambda_{k,p}| \gtrsim_p \sqrt{\frac{2^k}{k}}.
\]
\end{lemma}

The normalized Hadamard matrix $H\in \R^{2^k\times 2^k}$ diagonalizes $M$; after permuting
the columns of $H$, assume that the first $r$ diagonal entries of $H^\top M H$ are all
$\lambda_{k,p}$. Define the spectrum-truncated version of $M$ as
\[
	\widetilde{M} = H\diag\{\underbrace{\lambda_{k,p},\dots,\lambda_{k,p}}_{r\text{ copies}}, 0,\dots,0\}H^\top.
\]

The following lemma is a parameterized version of \cite[Lemma 3.6]{LWW21}, which follows the same proof.
\begin{lemma}\label{lem:decomposition}
	For every sufficiently small $\delta > 0$, there is a subset $I\subseteq U$ with $|I|\gtrsim \delta^2 r$ such that the $i$-th row of $\widetilde{M}$ ($i\in I$) can be decomposed as 
	\[
		\widetilde{M}_i = R_i + P_i,
	\]
	where each $R_i$ lies in $\range(\widetilde{M})$, $\{R_i\}_{i\in I}$ is a set of orthogonal vectors, $\norm{R_i}_2\geq \sqrt{1-\delta^2}\|\widetilde{M}_i\|_2$, $P_i$ is the orthogonal projection of $\widetilde{M}_i$ onto the subspace spanned by $\{R_j\}_{j\in I\setminus \{i\} }$ and $\norm{P_i}_2\leq \delta\|\widetilde{M}_i\|_2$.
\end{lemma}

Let $I$ be a set given by the preceding lemma with $\delta$ a sufficiently small absolute constant.
For every $i\in I$, choose independently a random sign matrix $S_i\in \{-1,1\}^{s\times s}$ and define $W_i = (1/\sqrt{s})S_i$. For each $j\in U$, define
\[
	X_j = \sum_{i\in I} \frac{(R_i)_j}{\norm{R_i}_2} W_i.
\]
For each $j\in U$, we also define
\[
	C_j = 
	\begin{pmatrix}
		I & 0 \\
		0 & I 
	\end{pmatrix}
	+
	\eta
	\begin{pmatrix}
		0 & X_j \\
		X_j^\top & 0
	\end{pmatrix}.
\]

Note that
\[
	Y_i := \sum_{j\in U} \widetilde{M}_{i, j} X_j = \sum_{i' \in I} \left\langle \widetilde{M}_i, \frac{R_{i'}}{\norm{R_{i'}}_2}\right\rangle W_{i'} = \norm{R_i}_2 W_i + E_i,
\]
where
\[
	E_i = \sum_{i'\in I\setminus\{i\}} \left\langle \widetilde{M}_i, \frac{R_{i'}}{\norm{R_{i'}}_2}\right\rangle W_{i'}.
\]

Define
\[
	\cS = \left\{ \{S_i\}_{i\in I}: \max_{j\in U} \norm{X_j}_{\op} \leq C_1 \text{ and } \max_{i\in I} \frac{\norm{E_i}_{\op}}{\|\widetilde{M}_i\|_2} \leq C_2\delta \right\}.
\]
\begin{lemma}\label{lem:S_size}
	There exist absolute constants $C_1, C_2 > 0$ such that $|\cS|\geq 2^{\Omega(|I|s^2)}$.
\end{lemma}

Suppose that $\eta$ is a sufficiently small absolute constant. By Lemma~\ref{lem:S_size}, for each such collection $\{S_i\}_{i\in I}$, we have $(3/4) I_{2s} \preceq C_j
\preceq (5/4)I_{2s}$ for all $j\in U$.

By \cite[Lemma 7.1]{LWW21} and the same net argument used in \cite[Lemma 7.2]{LWW21} after normalization there
exists
\[
	T\in\R^{m_0\times (2s+1)}, \qquad m_0 = O_p\mleft(\mleft(\frac{s\log(1/\eps)}{\eps^2}\mright)^{\max\{1,p/2\}}\mright)
\]
such that $\norm{Ty}_p^p = (1\pm\eps)\norm{y}_2^p$ for all $y\in\R^{2s+1}$, and with entries of magnitude at most $\poly(d/\eps)$. Write $T_\ell = (g_\ell, h_\ell)\in \R\oplus \R^{2s}$ for $\ell \in [m_0]$. Then we form a matrix $A$ of $2^k m_0$ rows, indexed by $(i,\ell)\in U\times [m_0]$, by	
\[
	A_{i,\ell} = (g_\ell i^\top, h_\ell^\top C_i^{1/2}) \in \R^{k}\oplus \R^{2s}.
\]
The hard matrix $\widetilde{A}$ is obtained from $A$ by the rounding described in Section~\ref{sec:rounding}.

\subsection{Analysis of Matrix $M$}
We prove Lemma~\ref{lem:eigenvalue} in this section.

\begin{proof}[Proof of Lemma~\ref{lem:eigenvalue}]
	Recall that $k=4h+1$ and $m=2h$. 
	The same argument as in \cite{LWW21} showed that there are $\binom{k}{m}$ eigenvalues of the same value and so we only prove the lower bound here. 
	
	Let $q = p-2$. We first assume that $p > 1$, so $q > -1$. Following the coefficient calculation in \cite[Lemma 3.2]{LWW21} with the odd-degree polynomial identity 
	\[
		(1 - x)^{2h} (1 + x)^{2h + 1} = (1 - x^2)^{2h} (1 + x)
	\]
	in place of $(1-x)^h (1+x)^h = (1-x^2)^{2h}$, we obtain that
	\begin{equation}\label{eqn:lambda_expression}
		\lambda_{k,p} = \sum_{i=0}^{2h} (-1)^i \binom{2h}{i} (|4h+1-4i|^q + |4h-1-4i|^q).
	\end{equation}
	Following the same argument in \cite{LWW21}, we have for $-1 < q < 2h$, 
	\[
		\lambda_{k,p} = (-1)^h c_q 2^k I_q,
	\]
	where
	\[
		c_q = -\frac{2\Gamma(q+1)}{\pi}\sin\frac{\pi q}{2},\qquad I_q = \int_0^\infty \frac{\sin^{2h}t \cos^{2h+1}t}{t^{q+1}}\, dt.
	\]
	It then suffices to lower bound the integral $I_q$. Note that
	\[
		I_q = \sum_{j=0}^\infty (-1)^j I_q^{(j)},
	\]
	where
	\[
		I_q^{(j)} = \int_0^{\pi/2} \left(\frac{1}{(j\pi+t)^{q+1}} - \frac{1}{((j+1)\pi+t)^{q+1}}\right) \sin^{2h}t \cos^{2h+1}t \, dt.
	\]
	Since $I_q^{(j)} > 0$ and $I_q^{(j)}$ is decreasing in $j$, we know that $I_q\geq I_q^{(0)} - I_q^{(1)}$. Write
	\[
		I_q^{(0)} - I_q^{(1)} = \int_0^{\pi/2} g(t) \sin^{2h}t \cos^{2h+1}t\, dt,
	\]
	where
	\[
		g(t) = \frac{1}{t^{q+1}} - \frac{1}{(\pi-t)^{q+1}} - \frac{1}{(\pi+t)^{q+1}} + \frac{1}{(2\pi-t)^{q+1}},\quad t\in \left[0,\frac{\pi}{2}\right].
	\]
	On the interval $[\pi/4-c_1/\sqrt{h}, \pi/4+c_1/\sqrt{h}]$, we have
	\[
		\sin^{2h}t \cos^{2h+1}t \geq \frac{c_2}{2^{2h}},\quad g(t)\geq c_{3,q}.
	\]
	Therefore,
	\[
		I_q \gtrsim_q \frac{2^{-k/2}}{\sqrt k}
	\]
	and so for $p>1$,
	\begin{equation}\label{eqn:lambda_lb}
		|\lambda_{k,p}| \gtrsim_p \frac{2^{k/2}}{\sqrt{k}}.
	\end{equation}
	
	Now we deal with $p=1$ (i.e.~$q=-1$), which follows from a limiting argument. Taking $p\searrow 1$ (i.e.~$q\searrow -1$) in \eqref{eqn:lambda_expression}, since it is a finite sum and is continuous in $q$, we have $\lambda_{k,p}\to \lambda_{k,1}$. One can show that $g(t)\geq c_4(q+1)$ for some absolute constant $c_4$ on the small neighbourhood around $\pi/4$, and it is easy to verify using the expression of $c_q$ that $c_q\asymp 1/(q+1)$. Hence, the hidden constant in \eqref{eqn:lambda_lb} can be made an absolute constant when $p$ is close to $1$. This proves the lemma for $p=1$.
\end{proof}

\subsection{Good Sign Matrices}

We prove Lemma~\ref{lem:S_size} in this section.

\begin{proof}[Proof of Lemma~\ref{lem:S_size}]
For brevity, let $c_i = R_i/\norm{R_i}_2$. Since $c_i$'s are orthonormal vectors, we know that $\sum_{i\in I} (c_i)_j^2\leq 1$ for every $j\in U$. By the non-commutative Khintchine inequality for the operator norm,
\[
	\Pr\left\{ \norm{X_j}_{\op} \geq t \right\} \leq 2e^{-ct^2 s},
\]
where $c>0$ is an absolute constant. Again, since $c_i$'s are orthonormal, 
\[
	\sum_{i'\in I\setminus\{i\}} \langle \widetilde{M}_i, c_{i'}\rangle^2 = \norm{P_i}_2^2 \leq \delta^2\norm{\widetilde{M}_i}_2^2
\]
and again by the non-commutative Khintchine inequality,
\[
	\Pr\left\{ \norm{E_i}_{\op}\geq t\delta\norm{\widetilde{M}_i}_2 \right\} \leq 2e^{-ct^2 s}.
\]
Taking a union bound over $j\in U$ and $i\in I$, and using $|I|\leq 2^k$ and $s=\Omega(k)$, we see that
\[
	\max_{j\in U} \norm{X_j}_{\op} \leq C_1\qquad\text{and}\qquad \max_{i\in I} \frac{\norm{E_i}_{\op}}{\norm{\widetilde{M}_i}_2} \leq C_2\delta
\]
for sufficiently large absolute constants $C_1, C_2 > 0$, except with probability at most
\[
	4\cdot 2^k e^{-c\min\{C_1^2,C_2^2\}s} < 1/10.
\]
This finishes the proof.
\end{proof}

\subsection{Rounding}\label{sec:rounding}

We use the same entrywise perturbation principle as in the proof of \cite[Theorem 5.1]{LWW21}. We first note that, for $x = (v, w) \in \R^k \oplus \R^{2s}$,
\[
	\norm{Ax}_p^p \geq c_p 2^k \norm{x}_2^p.
\]
Indeed, by the definition of $A$ and the embedding guarantee for $T$,
\[
	\norm{Ax}_p^p \geq \frac12 \sum_{j\in U}\left(\ip{j}{v}^2 + w^\top C_jw\right)^{p/2}
	\geq c_p\sum_{j\in U}\left(\ip{j}{v}^2+\norm{w}_2^2\right)^{p/2}.
\]
The last sum is at least $c_p2^k(\norm{v}_2^2+\norm{w}_2^2)^{p/2}$, since it is bounded below by both $2^k\norm{w}_2^p$ and $\sum_{j\in U}|\ip{j}{v}|^p\geq c_p2^k\norm{v}_2^p$ by the Khintchine lower bound.
Fix a sign matrix collection $S\in \cS$. Round every entry of $A$ to a common grid of spacing $\zeta$, resulting in $\widetilde{A}$. Since $A$ has $n = 2^k m_0$ rows,
\[
	\norm{(A-\widetilde{A})x}_p \leq n^{1/p}\zeta\sqrt{d}\norm{x}_2.
\]
Choosing
\[
	\zeta \leq c_p'\frac{\eps}{m_0^{1/p}\sqrt{d}},
\]
with $c_p'>0$ sufficiently small gives
\[
	\norm{(A-\widetilde{A})x}_p \leq C_p'\eps \norm{Ax}_p.
\]
Thus, after adjusting the constant in the choice of $\zeta$,
\[
	\norm{\widetilde{A}x}_p^p = (1\pm \eps)\norm{Ax}_p^p, \quad \forall x\in\R^d.
\]
Now we examine the magnitude of the entries of $A$. By the choice of $T$, every coordinate of $T_\ell=(g_\ell,h_\ell)$ has magnitude at most $\poly(d/\eps)$, and hence $\norm{h_\ell}_2\leq \poly(d/\eps)$. Since $\norm{C_j^{1/2}}_{\op}\leq \sqrt{5/4}$, every entry of $(g_\ell j^\top,h_\ell^\top C_j^{1/2})$ has magnitude at most $\poly(d/\eps)$. Since $\zeta^{-1}=\poly(d/\eps)$, the rounded entries use $O_p(\log(nd))$ bits.

\section{Recovery of Sign Matrices}

For $S=\{S_i\}_{i\in I}\in\cS$, $v\in \R^k$, $\tau \geq 0$ and
$w\in \Sph^{2s-1}$, define
\[
	F_S(v,\tau,w)
	:= \sum_{j\in U} (\langle j,v\rangle^2 + \tau w^\top C_j w)^{p/2}.
\]
All quantities $W_i,X_j,C_j,Y_i,E_i$ and $A$ are understood to be constructed
from the displayed collection $S$; when needed, we write $C_j^S$ or $A_S$ to
emphasize this dependence.
Fix $S\in\cS$ in the following display and in Lemma~\ref{lem:F_properties}, and let $A=A_S$ be the corresponding unrounded matrix.
When this $S$ is clear from the context, we write $F$ instead of $F_S$. Then
\begin{equation}\label{eqn:A_approximates_F}
	(1-\eps)F(v,\tau,w)
	\leq 
	\norm{ A \begin{pmatrix}
			v\\
			\sqrt{\tau}w
	\end{pmatrix} }_p^p 
	\leq (1+\eps)F(v,\tau,w).
\end{equation}
Hence, we shall first analyse some properties of the function $F$.

\begin{lemma}\label{lem:F_properties}
Suppose that $w\in \Sph^{2s-1}$ and $i\in I$ are arbitrary. The following properties hold.
\begin{enumerate}[label=(\roman*)]
	\item There exists an absolute constant $C_E > 0$ such that
	\[
		\left|	\frac{(2/p)\partial_\tau F(i,0,w) - b_0}{\eta\|R_i\|_2} 
		- w^\top \begin{pmatrix}
			        0 & W_i \\
			        W_i^\top & 0 \\
				 \end{pmatrix} w \right|
		\le C_E\delta,
	\]
	where $b_0 = (M\mathbf{1})_i$ is a constant independent of $i$.
	
	\item There is a constant $C_p > 0$ depending only on $p$ such that whenever $0\leq \tau \leq 1/4$,
	\[
		F(i,\tau,w)	\le C_p 2^k k^{p/2}.
	\]
\end{enumerate}
\end{lemma}
\begin{proof}
	Since $M$ and $\widetilde{M}$ agree on $\range(\widetilde{M})$, we have
	$\sum_{j\in U} M_{i,j} X_j = \sum_{j\in U}\widetilde{M}_{i,j}X_j = Y_i$. Therefore,
	\[
		\frac{\partial F}{\partial \tau}(i,0,w)
		= \frac{p}{2}\sum_{j\in U}|\langle i,j\rangle|^{p-2} w^\top C_j w 
		= \frac{p}{2}\left( \sum_{j\in U}|\langle i,j\rangle|^{p-2} + \eta w^\top \begin{pmatrix}
			0 & Y_i \\	Y_i^\top & 0 \\
		\end{pmatrix} w \right).
	\]
	Let $b_0 = \sum_{j\in U} |\ip{i}{j}|^{p-2}$, which is independent of $i$. We now have
	\[
		\frac{(2/p)\partial_\tau F(i,0,w) - b_0}{\eta\norm{R_i}_2} = w^\top \begin{pmatrix}
			0 & W_i \\	W_i^\top & 0 \\
		\end{pmatrix} w + \frac{1}{\norm{R_i}_2}
		w^\top \begin{pmatrix}
			0 & E_i \\	E_i^\top & 0 \\
		\end{pmatrix} w
	\]
	By the definition of $\cS$ and $\norm{R_i}_2\geq \sqrt{1-\delta^2}\norm{\widetilde{M}_i}_2$, we have
	\[
		\left| \frac{1}{\norm{R_i}_2}
			w^\top \begin{pmatrix}
				0 & E_i \\	E_i^\top & 0 \\
			\end{pmatrix} w \right| \leq C_E\delta.
	\]
	This proves (i).
	
	Next we prove (ii). Since $(3/4)I \preceq C_j\preceq (5/4)I$, we know that $w^\top C_j w \leq 5/4$. Hence
	\[
		F(i,\tau,w) \leq C_p\sum_{j\in U}\left(|\ip{i}{j}|^p+1\right).
	\]
	For uniformly random $j\in U$, $\ip{i}{j}$ has the distribution of a sum of $k$ independent Rademacher variables. Therefore $\E_j|\ip{i}{j}|^p\leq C_p k^{p/2}$ by the Khintchine inequality. It follows that
	\[
		F(i,\tau,w) \leq C_p2^k(k^{p/2}+1) \leq C_p2^k k^{p/2}. \qedhere
	\]
\end{proof}

Next we shall show how to recover sign matrices $\{S_i\}_{i\in I}$ from the quadratic form $w^\top \begin{psmallmatrix}
	0 & W_i \\ W_i^\top & 0\\
\end{psmallmatrix} w$, provided an accurate estimate of $\partial_\tau F_S(i,0,w)$, and then show how we can obtain such an accurate estimate.

Let $S = \{S_i\}_{i\in I}$ and $S' = \{S_i'\}_{i\in I}$ be two collections of sign matrices in $\cS$. We define their Hamming distance as
\[
	d_H(S,S') = \sum_{i\in I} d_H(S_i, S_i').
\]

\begin{lemma}\label{lem:conditional_recovery}
	There exists an absolute constant $c_0>0$ such that the following holds. Let $\cC \subseteq \cS$ be a family of sign matrix collections with pairwise Hamming distance at least $c_0 |I| s^2$. For each $i\in I$, there is a function $\hat{D}_i:\Sph^{2s-1}\to \R$. We say $S\in\cC$ is \emph{compatible} with $\{\hat{D}_i\}_{i\in I}$ if 	
	\begin{equation}\label{eqn:D_hat}
		\left|\hat{D}_i(w) - \frac{\partial F_S}{\partial \tau}(i,0,w)\right| \leq \theta\eta\norm{R_i}_2
	\end{equation}
	for all $i\in I$ and $w\in\Sph^{2s-1}$. If	
	\begin{equation}\label{eqn:important_condition}
		\frac{2\theta}{p} + C_E\delta \leq \frac{1}{2}\sqrt{c_0},
	\end{equation}
	then at most one $S\in \cC$ is compatible with $\{\hat{D}_i\}_{i\in I}$.
\end{lemma}
\begin{proof}
	Suppose that $d_H(S,S')\geq c_0|I|s^2$, then there exists $i\in I$ such that  $d_H(S_i,S_i')\geq c_0 s^2$. Let $W_i = S_i/\sqrt{s}$ and $W_i' = S_i'/\sqrt{s}$, then
	\[
		\norm{W_i - W_i'}_F^2 = \frac{4}{s}d_H(S_i,S_i') \geq 4 c_0 s.
	\]
	Hence
	\[
		\norm{W_i - W_i'}_{\op}\geq \frac{1}{\sqrt s}\norm{W_i - W_i'}_F \geq 2\sqrt{c_0}
	\]
	and there exists $w\in \Sph^{2s-1}$ such that
	\begin{equation}\label{eqn:direction_w}
		\left|w^\top \begin{pmatrix}
			0 & W_i \\
			W_i^\top & 0 \\
		\end{pmatrix} w - w^\top \begin{pmatrix}
		0 & W_i' \\
		W_i'^\top & 0 \\
	\end{pmatrix} w\right| \geq 2\sqrt{c_0}.
	\end{equation}
	On the other hand, by Lemma~\ref{lem:F_properties}(i), we have
	\[
		\left| \frac{(2/p)\hat{D}_i(w) - b_0}{\eta\norm{R_i}_2} - w^\top \begin{pmatrix}
			0 & W_i \\
			W_i^\top & 0 \\
		\end{pmatrix} w \right| \leq \frac{2\theta}{p} + C_E\delta \leq \frac{1}{2}\sqrt{c_0}.
	\]
	If $\{\hat{D}_i\}_{i\in I}$ is compatible with both $S$ and $S'$, then by the triangle inequality, the left-hand side of \eqref{eqn:direction_w} would be at most $\sqrt{c_0}$, which contradicts \eqref{eqn:direction_w}.	
\end{proof}

\begin{lemma}\label{lem:derivative_approximation}
	Suppose that an instance of $Q_p$ is a correct $\ell_p$-subspace sketch for the rounded hard matrix $\widetilde{A}$.
	Then one can construct functions $\hat{D}_i:\Sph^{2s-1}\to\R$ for each $i\in I$ such that \eqref{eqn:D_hat} holds for all $w\in \Sph^{2s-1}$.
\end{lemma}
\begin{proof}
	We claim that there exist an integer $L\asymp \log(1/\eps)$ and coefficients $\beta_0,\dots,\beta_L$ such that $\sum_\ell |\beta_\ell| \leq C L^2$ and
	\[
		\left| \sum_{\ell=0}^L \beta_\ell F_S(i,\tau_\ell,w) - \partial_\tau F_S(i,0,w) \right| \leq C_p \eps L^2 2^k k^{p/2}.
	\]
	The proof of the claim is postponed to Section~\ref{sec:diff_interpolation}.
	
	Let $x_\ell = (i, \sqrt{\tau_\ell}w)$. By the correctness of $Q_p$ for $\widetilde{A}$, the rounding guarantee from Section~\ref{sec:rounding}, and \eqref{eqn:A_approximates_F}, we have
	\[
		\left| Q_p(x_\ell) - F_S(i,\tau_\ell,w) \right| \leq C_p\eps F_S(i,\tau_\ell,w).
	\]
	By Lemma~\ref{lem:F_properties}(ii), we have
	\[
		\left| \sum_{\ell=0}^L \beta_\ell  Q_p(x_\ell) - \sum_{\ell=0}^L \beta_\ell F_S(i,\tau_\ell,w) \right| \leq C_p \eps 2^k k^{p/2} \sum_{\ell=0}^L |\beta_\ell| \leq C_p \eps L^2 2^k k^{p/2},
	\]
	where the absolute constant in $\sum_\ell|\beta_\ell|\leq C L^2$ has been absorbed into $C_p$.
	Define $\hat{D}_i(w) = \sum_{\ell=0}^L \beta_\ell Q_p((i, \sqrt{\tau_\ell}w))$, we have
	\[
		\left|\hat{D}_i(w) - \partial_\tau F_S(i,0,w)  \right| \leq C_p \eps L^2 2^k k^{p/2}.
	\]
	Note that 
	\[
		\norm{R_i}_2 \geq \sqrt{1-\delta^2}\norm{\widetilde{M}_i}_2 = \sqrt{1-\delta^2}|\lambda_{k,p}|\sqrt{\frac{r}{2^k}} \gtrsim_p \sqrt{\frac{r}{k}} \gtrsim \frac{2^{k/2}}{k^{3/4}}.
	\]
	Recall the assumption \eqref{eqn:assumption_k}, we have, after adjusting constants, $C_p \eps L^2 2^k k^{p/2} \leq \theta \eta \norm{R_i}_2$.
\end{proof}

\subsection{Interpolation of Derivatives}\label{sec:diff_interpolation}
\begin{lemma}
	Let $\gamma > 0$. There are constants $R > 1$ and $C_\gamma > 0$ such that, for every $N\geq 1$, there exist nodes
	\[
		0 = t_0 < t_1 < \cdots < t_N = \frac{1}{4}
	\]
	and coefficients $\beta_0,\beta_1,\dots,\beta_N$ for which
	\[
		\left| \sum_{i=0}^N \beta_i (a + ct_i)^\gamma - \gamma c a^{\gamma-1} \right| \leq C_\gamma N^2 R^{-N} a^\gamma
	\]
	whenever $a\geq 1$ and $3/4 \leq c \leq 5/4$. Furthermore, the coefficients can be chosen such that
	\[
		\sum_{i=0}^N |\beta_i| \leq C N^2.
	\]
\end{lemma}
\begin{proof}
	Define $f(t) = (a + ct)^\gamma$, then $f'(0) = \gamma c a^{\gamma-1}$. Let $p_N$ be the Lagrange interpolating polynomial of degree at most $N$ interpolating $f$ at the nodes $t_i = (1/8)(1-\cos((i/N)\pi))$. We know that $p_N'(0) = \sum_{i=0}^N \beta_i f(t_i)$, where
	\[
		-\frac{1}{8}\beta_i = 
		\begin{cases}
			(2N^2+1)/6 & i = 0\\
			2(-1)^i/(1-\cos(i\pi/N)) & 1\leq i\leq N-1 \\
			(-1)^N/2 & i=N.
		\end{cases}
	\]
	It is easy to verify that $\sum_i |\beta_i| \leq CN^2$. It remains to bound $|p_N'(0) - f'(0)|$. 
	
	Suppose that $t(x)$ is a linear function that maps $[-1,1]$ to $[0,1/4]$. Let $g(x) = f(t(x))$ and $t(x_i) = t_i$ for $x_i = \cos((N-i)\pi/N)$, which are exactly Chebyshev nodes. We have $g'(-1) = 1/8 f'(0)$. Also let $w_N(x) = p_N(t(x))$, then $w_N'(-1) = p_N'(0)/8$. Thus
	\[
		|p_N'(0) - f'(0)| = 8|w_N'(-1) - g'(-1)|.
	\]
	 By \cite[Corollary 4.6]{tadmor}, 
	 \[
	 	|w_N'(-1) - g'(-1)| \leq C_\alpha \cdot \max_{z\in E_\alpha} |g(z)| \cdot N^{\frac{5}{2}} \alpha^{-N},
	 \]
	 where $\alpha$ is chosen such that $g$ is analytic in the Bernstein ellipse $E_\alpha$. It is easy to see that we can take $\alpha = 3/2$ and $|g(z)|\leq (3a/2)^\gamma$. Hence,
	 \[
	 	|w_N'(-1) - g'(-1)| \leq C a^\gamma N^{\frac{5}{2}} \left(\frac{2}{3}\right)^N.
	 \]
	 Choose $R = 4/3$, then
	 \[
	 	N^{\frac{5}{2}} \left(\frac{2}{3}\right)^N = N^2 \left(\frac{3}{4}\right)^N \left( N^{\frac{1}{2}} \left(\frac{8}{9}\right)^N \right) \leq C N^2 \left(\frac{3}{4}\right)^N = C N^2 R^{-N}.
	 \]
	 This completes the proof.
\end{proof}

The claim in the proof of Lemma~\ref{lem:derivative_approximation} follows from applying the preceding lemma with $\gamma = p/2$, $N \asymp \log(1/\eps)$ such that $R^{-N} \leq \eps$ and Lemma~\ref{lem:F_properties}(ii).

\subsection{Lower Bound for Subspace Sketches}
Let $C_E$ denote the absolute constant from Lemma~\ref{lem:F_properties}, and let
$c_0>0$ be a sufficiently small absolute constant for Lemma~\ref{lem:conditional_recovery}
and the packing argument below.
The absolute constant $\theta>0$ in the construction is chosen small enough that
$2\theta/p\leq \frac{1}{4}\sqrt{c_0}$, and then choose $\delta>0$ small enough
that \eqref{eqn:important_condition} holds.
By Lemma~\ref{lem:S_size} and a greedy Gilbert-Varshamov packing inside $\cS$, there is a subfamily $\cC\subseteq\cS$ with pairwise Hamming distance at least $c_0|I|s^2$ and $|\cC|=2^{\Omega(|I|s^2)}$. 

By Yao's minimax principle, it suffices to consider a deterministic data structure $Q_p$ for the distribution that chooses $S$ uniformly from $\cC$ and gives the rounded hard matrix $\widetilde{A}_S$ as input. Suppose $Q_p$ is correct with probability at least $0.9$ over this choice of $S$, and let $\cC'\subseteq\cC$ be the set of $S$ for which $Q_p$ is correct. Then $|\cC'|\geq 0.9|\cC|$, and $\cC'$ still has pairwise Hamming distance at least $c_0|I|s^2$. Combining Lemmata~\ref{lem:conditional_recovery} and~\ref{lem:derivative_approximation}, we see that, for each $S\in\cC'$, the deterministic sketch $Q_p$ can recover $S$. Applying the deterministic counting argument to $\cC'$, the sketch must have at least
\[
	\log|\cC'| = \Omega(|I|s^2)
	= \Omega\mleft(\frac{d^2}{\eps^2\polylog(1/\eps)}\mright)
\]
bits. Finally, the constructed matrix has $n_0=2^k m_0$ rows. Since $2^k=\widetilde\Theta_p(\eps^{-2})$, $s=\Theta(d)$ and $m_0=O_p((s\log(1/\eps)/\eps^2)^{\alpha_p})$, we have
\[
	n_0=\widetilde O_p(d^{\alpha_p}/\eps^{2+2\alpha_p}).
\]
For any larger number of rows $n$, we pad the hard instance with zero rows. This proves Theorem~\ref{thm:subspace_sketch_lower_bound}.

\section{Dimension Lower Bound}
\begin{proof}[Proof of Theorem~\ref{thm:dimension_lower_bound}]
	We use the standard conversion from sketch lower bounds to dimension lower bounds
	from \cite{LWW21}. Let $\widetilde A$ be the rounded hard instance from
	Theorem~\ref{thm:subspace_sketch_lower_bound}, and view its column span as a
	$d$-dimensional subspace of $\ell_p^{n_0}$; indeed, $\widetilde A$ has full column rank by the conditioning and rounding bounds in Section~\ref{sec:rounding}. If this
	subspace admitted a $(1+\delta)$-embedding into $\ell_p^N$, with
	$\delta = c_p\eps$ sufficiently small, then there would be a matrix
	$B\in\R^{N\times d}$ such that
	\[
		\norm{Bx}_p^p=(1\pm \eps)\norm{\widetilde A x}_p^p
		\qquad\text{for all }x\in\R^d .
	\]
	Rounding the entries of $B$ as in \cite{LWW21} gives a valid
	$\ell_p$-subspace sketch for $\widetilde A$ using
	$O_p(Nd\log(n_0d/\eps))$ bits. Therefore Theorem~\ref{thm:subspace_sketch_lower_bound}
	implies
	\[
		N d \log(n_0 d/\eps)
		\gtrsim_p
		\frac{d^2}{\eps^2\polylog(1/\eps)}.
	\]
	Since $n_0=\widetilde O_p(d^{\alpha_p}/\eps^{2+2\alpha_p})$, rearranging gives
	\[
		N \gtrsim_p \frac{d}{\eps^2\polylog(d/\eps)} .
	\]
\end{proof}

\section*{Acknowledgements}
The author is supported in part by the Singapore Ministry of Education AcRF Tier 1 grant RG21/25.

The central technical idea of this paper originated from ChatGPT 5.6 Sol. The author substantially revised and completed the AI-assisted material, and is responsible for all mathematical claims, proofs, references, and the final manuscript.

\bibliographystyle{plain}
\bibliography{reference}

\end{document}
